\documentclass[11pt]{article}
\usepackage[a4paper,margin=27mm]{geometry}
\usepackage{amsmath,amssymb,mathtools}
\usepackage{fontspec,unicode-math}
\setmainfont{Latin Modern Roman}
\setmathfont{Latin Modern Math}
\usepackage{xurl,hyperref}
\hypersetup{colorlinks=true,urlcolor=blue,linkcolor=blue}
\newcommand{\sref}[2]{\hyperlink{source:#1:#2}{[S#2]}}
\newcommand{\cataloguescope}[1]{\paragraph{Mathematical question.} #1}
\begin{document}
\section{Conjecture 4.12}\label{statement}

\subsection{Definitions and mathematical statement}

\paragraph{Shared notation.}$\mathbb N=\{1,2,\ldots\}$, and $\mathbb Z,\mathbb Q,\mathbb R,\mathbb C$ denote the integers, rationals, reals and complex numbers. Any other conventions are specified locally.


For an algebraic $\gamma$, define $h(\gamma)=\deg(\gamma)^{-1}\log M(\gamma)$ using its primitive minimal-polynomial Mahler measure. For an algebraic tuple $\gamma=(\gamma_1,\ldots,\gamma_m)$ put $d(\gamma)=[\mathbb Q(\gamma):\mathbb Q]$ and $\mu(\gamma)=d(\gamma)\max_i h(\gamma_i)$. Fix $\theta\in\mathbb C^m$ with $t=\operatorname{trdeg}_{\mathbb Q}\mathbb Q(\theta)\ge1$.
\paragraph{Statement recorded in the supplied source.}
Do there exist $c_1,c_2>0$, depending on $\theta$, such that for every pair of real sequences $(D_\nu),(\mu_\nu)$ ($\nu\ge0$) satisfying $c_1\le D_\nu\le\mu_\nu$, $D_\nu\le D_{\nu+1}\le2D_\nu$, $\mu_\nu\le\mu_{\nu+1}\le2\mu_\nu$, and unbounded $\mu_\nu$, infinitely many $\nu$ admit an algebraic tuple $\gamma$ with \[d(\gamma)\le D_\nu,\qquad \mu(\gamma)\le\mu_\nu,\qquad \max_i|\theta_i-\gamma_i|\le\exp(-c_2D_\nu^{1/t}\mu_\nu)?\] \sref{statement}{2}
\subsection{Short English statement}
Can every tuple of positive transcendence degree be simultaneously approximated at the stated sharp rate along infinitely many indices of every controlled degree-and-height growth schedule?
\subsection{Sources}
\begin{description}
\item[\hypertarget{source:statement:1}{S1}] ULAM.AI and the UnsolvedMath Contributors, UnsolvedMath: A Curated Collection of Open Mathematics Problems, catalogue entry. CC BY 4.0. \url{https://huggingface.co/datasets/ulamai/UnsolvedMath}.
\item[\hypertarget{source:statement:2}{S2}] Michel Waldschmidt, Open Diophantine Problems, arXiv:math/0312440v2 (24 January 2004), Conjecture 4.12 and preceding definition of mu. \url{https://arxiv.org/abs/math/0312440}.
\end{description}
\label{end:statement}
\end{document}
